\section{Academic Integrity in the Age of AI}\label{sec:integrity}

The approach is \emph{clarity, design, evidence, and fairness}: prevent misuse by making expectations clear and assessments robust, and handle suspected misuse proportionately.

\subsection{Course-level AI-use levels}
Each faculty member declares one of five levels in the syllabus (and may vary it by assessment). The scale is an original VSU proposal; its graded-levels idea resembles published AI assessment scales \citep{perkins2024}, but the levels here differ and attribution is to be confirmed \todoverify{compare with source before claiming adaptation}.

\begin{table}[ht]
\centering
\caption{Proposed AI-use levels.}
\small
\begin{tabular}{@{}cp{3.0cm}p{5.8cm}p{4.2cm}@{}}
\toprule
Level & Name & What students may do & Disclosure \\
\midrule
0 & No AI & Work completed without AI (e.g., in-class, oral, or supervised tasks). & None needed; supervised conditions. \\
1 & Study support & AI for understanding, practice, and planning only; no AI-generated text in submitted work. & None, unless asked. \\
2 & AI-assisted (default) & AI for brainstorming, feedback, editing, and language support; student authors the substance. & Short statement. \\
3 & AI-collaborative & AI may generate drafts or code that students critique, verify, and revise. & Statement plus process evidence. \\
4 & AI-integrated & AI use is expected and assessed as part of the task. & Full documentation of prompts or workflow. \\
\bottomrule
\end{tabular}
\end{table}

If a course is silent, \textbf{Level~2} applies. This default is proposed because it permits supportive use while keeping authorship with the student, avoiding both a blanket ban and unrestricted generation; the panel is asked to confirm it. Faculty may use the declaration templates in Appendix~\ref{app:declarations} and may set different levels for different assessments. Mixed-level courses should state the level beside each assessment.

\subsection{Assessment redesign}
Generic written-output tasks are easy for AI to substitute for student work [CITATION NEEDED], so faculty are encouraged to:
\begin{itemize}
  \item match the AI-use level to the learning outcome being assessed (some outcomes, such as foundational skills, merit Level~0 or~1 conditions);
  \item use authentic, local, and context-specific tasks (e.g., Visayas community cases, field data, course discussions);
  \item add short oral defenses, in-class checkpoints, or reflections on process and decisions for higher-stakes work;
  \item stage large assignments (proposal, outline, draft, final) so that development is visible;
  \item combine lower-stakes open-AI work with a smaller number of supervised assessments; and
  \item test their own assessments with an AI tool to see what it can and cannot do.
\end{itemize}

\subsection{Process-based evidence}
Evidence of authorship is gathered through ordinary teaching activities rather than surveillance: version histories, drafts, annotated sources, short reflection notes, and brief conversations about the work. Students may be asked to explain their work; inability to explain submitted work is a signal for conversation, not proof of misuse.

\subsection{Cautions about AI-text detectors}
Published evaluations report that AI-text detection tools can misclassify text and can be biased against non-native English writers \citep{liang2023,weberwulff2023} [CITATION NEEDED: confirm findings]. Therefore:
\begin{itemize}
  \item detector output \textbf{alone shall never} be the basis for a finding of misconduct or a grade penalty;
  \item detectors are not mandated and the University does not endorse their routine use; and
  \item any use must be treated as, at most, one weak signal alongside process evidence and a conversation with the student.
\end{itemize}

\subsection{Proportionate and fair handling of suspected misuse}
\begin{enumerate}
  \item \textbf{Conversation first.} The instructor discusses the concern with the student and reviews the process evidence. Cases of this kind are intended to be resolved here as educational matters (e.g., unclear declaration, first minor lapse).
  \item \textbf{Educational outcomes for minor or first-time cases:} revision or resubmission, reduced marks for the affected component, or a learning agreement.
  \item \textbf{Referral for serious or repeated cases} to the existing student discipline or academic integrity process, with written notice of the concern and evidence.
  \item \textbf{Due process.} The student receives clear notice, access to the evidence, an opportunity to respond, and a decision by someone other than the original complainant for contested cases, plus an avenue of appeal consistent with existing University procedures. The standard of evidence is the one in existing \vsu{} rules. \todoverify{Align with the Student Handbook and Code of Conduct provisions.}
  \item \textbf{Fairness safeguards:} consideration of language background, disability, and the clarity of the declaration in force; no presumption of guilt from use of AI in itself where the declaration permitted it.
  \item \textbf{Records and learning:} anonymized case data inform training and assessment redesign.
\end{enumerate}

\subsection{Misuse that is clearly prohibited}
Regardless of level: submitting AI output as one's own work where the declaration prohibits it or requires disclosure and none is made; fabricating data or sources; impersonation in supervised assessments; uploading others' work or confidential materials into AI tools without permission; and using AI to harass, deceive, or create deepfakes of others.
